


% To place above and left with different values:
% https://tex.stackexchange.com/questions/16231/node-below-and-left-of-another-node-in-tikz#:~:text=Building%20on%20%40Seamus,end%7Bdocument%7D
% Example: \node[latentnode, above left= \verticalnodedistance*\nodedistance and \horizontalnodedistance*\nodedistance of s0] (s11) {};
%=============================================
% pgfplots
%---------------------------------------------
\usepackage{pgfplots}
\pgfplotsset{compat=newest}
\tikzstyle{every picture}+=[remember picture]
\tikzstyle{na} = [baseline=-.5ex]
\everymath{\displaystyle}
\usetikzlibrary{arrows,shapes}
\usetikzlibrary{positioning}
\usetikzlibrary{matrix}

\usetikzlibrary{calc}
\usetikzlibrary{patterns.meta}
\usetikzlibrary{arrows,fit,positioning,shapes}  
\usetikzlibrary{backgrounds}
\usetikzlibrary{decorations.pathreplacing}
\usetikzlibrary{bayesnet}
\usetikzlibrary{shapes.misc,shapes.geometric}
\usetikzlibrary{fit,backgrounds}
\usetikzlibrary{calligraphy}



%=============================================
% Create layers
%---------------------------------------------
\pgfdeclarelayer{back}
\pgfdeclarelayer{front}
\pgfsetlayers{back,main,front}
\makeatletter
\pgfkeys{%
  /tikz/on layer/.code={
    \pgfonlayer{#1}\begingroup
    \aftergroup\endpgfonlayer
    \aftergroup\endgroup
  },
  /tikz/node on layer/.code={
    \gdef\node@@on@layer{%
      \setbox\tikz@tempbox=\hbox\bgroup\pgfonlayer{#1}\unhbox\tikz@tempbox\endpgfonlayer\egroup}
    \aftergroup\node@on@layer
  },
  /tikz/end node on layer/.code={
    \endpgfonlayer\endgroup\endgroup
  }
}
\def\node@on@layer{\aftergroup\node@@on@layer}



%=============================================
% Define block styles
%---------------------------------------------
\tikzset{
  latentnode/.style  ={draw,minimum width=3em, shape=circle,thick, black, fill=white, anchor=base, text height=2ex, text depth=0.75ex, },
  line/.style={draw, -latex'},
  line2/.style={draw, -latex', dotted},
  jump left/.style={draw=none, inner xsep=0pt},
  jump line/.style={line, shorten <=5pt}
}

\tikzset{fontscale/.style = {font=\relsize{#1}}
    }



%=============================================
% Add symbols
%---------------------------------------------
\def\checkmark{\tikz\fill[scale=0.4](0,.35) -- (.25,0) -- (1,.7) -- (.25,.15) -- cycle;} % Checkmark

%=============================================
% Define layers
%---------------------------------------------
\pgfdeclarelayer{foreground} 
\pgfdeclarelayer{background}
   \pgfsetlayers{background,%
                 main,%
                 foreground%
                 }
                 
% Create table where it is written at 45° 
% https://tex.stackexchange.com/questions/78938/tables-in-latex-with-diagonal-first-row
\newcommand{\tikzmark}[1]{\tikz[remember picture,overlay]\coordinate (#1);}
\newcolumntype{T}[1]{@{\hspace{\tabcolsep}}c@{\hspace{\tabcolsep}\tikzmark{#1}}}
\newcolumntype{L}[1]{@{\hspace{\tabcolsep}}l@{\hspace{\tabcolsep}\tikzmark{#1}}}


%=============================================
% Pgfplot
%---------------------------------------------
\usepgfplotslibrary{fillbetween} % confidence interval, https://tex.stackexchange.com/questions/596613/how-to-add-confidence-regions-around-a-graph-from-a-table
\usepgfplotslibrary{groupplots} % several plots in an image https://tex.stackexchange.com/questions/243976/get-caption-subcaption-every-plot-using-groupplot-pgf-tikz

